For twenty editions, ICWSM has studied social life online as platforms and methods changed. AI systems now rank, rewrite, generate, and moderate those interactions. We examine what ICWSM's record carries forward into this shift. We analyze 2,139 indexed contributions from 2007 to 2026, separating topics, estimated with nonnegative matrix factorization, from problem framings, approximated with explicit textual cues. Platform contexts turn over: blog mentions fall, Twitter mentions rise and later decline, while Reddit and language-model research grow. Community research keeps a similar topic share (10.5\% to 10.0\%), while governance cues within that research rise from 4.0\% to 34.3\%. This contrast persists after removing governance-matching features and under fixed-length checks. Harm-related cues also increase when abstracts are restricted to a fixed length. Selected studies repeatedly confront the same methodological problems: whose behavior is observed, what a trace measures, which comparison supports a causal claim, whose judgment defines governance, and what data access makes observable. We show how AI-mediated interaction sharpens these problems and provide a reporting checklist for studying them across changing systems.
